%% RESUMO EM LINGUA ESTRANGEIRA (elemento obrigatorio -- NBR 6028)
%% Traducao fiel do resumo em portugues. As keywords sao definidas em
%% config/dados.tex.

\begin{abstract}
This research investigates the application of artificial intelligence to the
triage of repetitive court notices in public legal advocacy, taking as a case
study the Judgment Enforcement Unit (NCS) of the Attorney General's Office of
the Federal District (PGDF), a unit specialized in individual enforcement
proceedings arising from collective judgments. It starts from the observation
that, in repetitive and low-complexity cases, deciding on the appropriate
measure tends to be quick for an experienced professional, so that the
operational cost lies in the individual handling of each occurrence. In this
workflow, artificial intelligence can support the classification of the notice
and the suggestion of appropriate measures, while automation mechanisms can
take over the resulting operational steps. The theoretical framework combines
the literature on digital government and public sector innovation, the
transition from street-level to system-level bureaucracy, and the governance
of artificial intelligence in Public Administration, with attention to human
oversight, transparency and personal data protection. The research is applied
and uses a mixed-methods approach, carried out through a single case study
selected by the criterion of typicality. The compliance of the solution is
examined in light of the regulatory framework applicable to the PGDF, made up
of the General Data Protection Law, PGDF Ordinance No. 160/2025 and Federal
District Decree No. 48,901/2026. The research aims to offer a methodological
framework that other attorney general's offices could potentially replicate,
while keeping human responsibility for decision-making and for the performance
of procedural acts.
\end{abstract}
